% [原文 p.221]

\begin{proof}
首先归约到 $f^{-1}(0)\cup D\cup\Delta$ 为 snc 除子的情形。设 $p\colon S''\to S$ 是 $(S,f^{-1}(0)\cup D\cup\Delta)$ 的对数奇点消解，$D'':=p^{-1}_*D$，$s''\in D''\cap p^{-1}(s)$，$\Delta''$ 由
\[
K_{S''}+D''+\Delta''=p^*(K_S+D+\Delta)\equiv0
\]
定义。特别地，$(S,D+\Delta)$ 为 lc 当且仅当 $(S'',D''+\Delta'')$ 为 lc。又，$(D'',\Delta''|_{D''})$ 在 $s''$ 处为 lc 可由 $(D,\Delta|_D)$ 在 $s$ 处为 lc 推出。以上表明，只需讨论 $S''$，归约便完成了。

记 $\Delta=\sum_i a_iE_i$，如下定义 $\Delta_1$，则 $\Delta_1$ 有效。
\[
\Delta_1:=\sum_{E_i\subset f^{-1}(0)}E_i+\sum_{E_i\not\subset f^{-1}(0)}\min\{a_i,1\}E_i.
\]
对 $\Q$-分解且 dlt 的 $(S,D+\Delta_1)$ 运行 $T$ 上的 $(K+D+\Delta_1)$-MMP。由定理 3.47，各步都可执行，纲领终于下述曲面 $S'$。
\[
f\colon S\overset{g}{\dashrightarrow}S'\overset{h}{\to}T.
\]
记 $D':=g_*D$，$\Delta'_1:=g_*\Delta_1$，$\Delta'=g_*\Delta$。此时，$(S',D'+\Delta'_1)$ 是 $\Q$-分解的且为 dlt（推论 3.44），$K_{S'}+D'+\Delta'_1$ 为 $h$-nef。又，与上面同样，$(S,D+\Delta)$ 为 lc 当且仅当 $(S',D'+\Delta')$ 为 lc（引理 2.30）。

由 $K_{S'}+D'+\Delta'\equiv0$ 可知，$\Delta'_1-\Delta'\equiv K_{S'}+D'+\Delta'_1$ 为 $h$-nef。又由 $\Delta_1$ 的定义，$h_*(\Delta'-\Delta'_1)$ 有效。因此，或 $\Delta'=\Delta'_1$，或 $\Delta'-\Delta'_1$ 有效且 $\Supp(\Delta'-\Delta'_1)\subset h^{-1}(0)$（推论 3.39）。

由 $\Delta'-\Delta'_1\ge0$ 与 $\Delta_1$ 的定义，$\Supp\lfloor\Delta'\rfloor=\Supp\lfloor\Delta'_1\rfloor$。对 $K_{S'}+D'+\Delta'$ 应用推论 5.49，得到 $(\lfloor\Delta'_1\rfloor\cup D')\cap h^{-1}(0)$ 的连通性。

首先假定 $g(s)\notin\lfloor\Delta'_1\rfloor$。由连通性，$\lfloor\Delta'_1\rfloor=0$。由 $\Delta_1$ 的定义可知 $\Delta'=\Delta'_1$。因此 $(S',D'+\Delta')$、$(S,D+\Delta)$ 为 lc。

以下可假定 $g(s)\in\lfloor\Delta'_1\rfloor$。也就是说，$D'+\Delta'_1$ 含 $2$ 个包含 $g(s)$ 的既约成分。因此，在 $g(s)$ 的近旁 $\Delta'_1=\lfloor\Delta'_1\rfloor$，$D'$ 与 $S'$ 非奇异（推论 5.55）。又由 $K_D+\Delta|_D=g^*(K_{D'}+\Delta'|_{D'})$，$(D',\Delta'|_{D'})$ 在 $g(s)$ 处为 lc。这在 $g(s)$ 的近旁意味着 $\Delta'=\Delta'_1$（推论 5.55）。进一步，在整个 $S'$ 上 $\Delta'=$
% [原文 p.222]
$\Delta'_1$ 成立。于是 $(S',D'+\Delta')$、$(S,D+\Delta)$ 为 lc。
\end{proof}

% [Tutor Note: 命题 5.58 的证明用 $T$ 上的 MMP 代替直接消解，把问题推到曲面 $S'$，再靠推论 5.49 的连通性与推论 5.55 的二维分类处理两种情形。这一思路在伴随反转的归纳中反复出现。]
\par\medskip\noindent\textbf{注意}\quad 本证明与 [Sho92] 中定理 5.50 的证明类似。命题 5.58 在 $S$ 为（不必是曲面的）正规簇时也成立，只要 $\Q$-分解 dlt 版的 MMP 在 $\dim T$ 维时可执行。特别地，三维情形也成立（[K+92, 6-8]、[Sho92]）。又，证明只需稍作改动便几乎照原样可行。不过，条件 (2) 须改为“$(D,\Delta|_D)$ 在 $D\cap h^{-1}(0)$ 的某连通成分的近旁为 lc”。

伴随对 dlt 对也成立。

\begin{propositionp}{5.59}
设 $X$ 是正规簇，$S$ 是余维 $2$ 且为 Cartier 的既约 Weil 除子，$B=\sum b_iB_i$ 是有效 $\Q$-除子。若 $(X,S+B)$ 为 dlt，则 $(S,B|_S)$ 也为 dlt，且 $\lfloor B|_S\rfloor=\lfloor B\rfloor|_S$。
\end{propositionp}

\begin{proof}
由推论 5.52，$S$ 正规。$(X,S+B)$ 为 dlt，故存在闭部分概形 $Z\subset X$，使 $X\setminus Z$ 非奇异、$S+B|_{X\setminus Z}$ 为 snc 除子，且 $\operatorname{discrep}(\operatorname{center}\subset Z,X,S+B)>-1$ 成立。令 $Z':=Z\cap S$。此时 $S\setminus Z'$ 非奇异，$B|_{S\setminus Z'}$ 为 snc 除子，且
\[
\operatorname{discrep}(\operatorname{center}\subset Z',S,B|_S)\ge\operatorname{discrep}(\operatorname{center}\subset Z,X,S+B)>-1
\]
成立（命题 5.46）。关于最后的等式，先看 $\lfloor B|_S\rfloor\ge\lfloor B\rfloor|_S$，这是显然的。把推论 5.53 用于 dlt 对 $(X\setminus\Supp\lfloor B\rfloor,S+\{B\})$，得 $\Supp\lfloor B|_S\rfloor\subset\Supp\lfloor B\rfloor|_S$。再把推论 5.53 与命题 2.43 结合，得 $\lfloor(1-\epsilon)B|_S\rfloor=0$（$0<\epsilon\ll1$）。故 $\lfloor B|_S\rfloor=\lfloor B\rfloor|_S$。
\end{proof}

\begin{remark}{5.60}
命题 5.59 的逆不成立。注意 5.54 所给例子即成反例。
\end{remark}

\begin{propositionp}{5.61}
设 $X$ 是正规簇，$S$ 是余维 $2$ 且为 Cartier 的既约 Weil 除子，$B=\sum b_iB_i$ 是有效 $\Q$-除子。若 $(X,S+B)$ 为 dlt 且 $\lfloor B\rfloor$ 的各连通成分既约，则 $(S,B|_S)$ 为 plt。
\end{propositionp}

\begin{proof}
由命题 5.59，$(S,B|_S)$ 为 dlt。若它不为 plt，则由命题 5.51，存在 $\lfloor B|_S\rfloor$ 的相异既约成分 $D_1,D_2$，使 $\operatorname{discrep}(\operatorname{center}\subset D_1\cap D_2,S,B|_S)=-1$。由命题 5.46，$\operatorname{discrep}(\operatorname{center}\subset D_1\cap D_2,X,S+B)=-1$。由 $\lfloor B|_S\rfloor=\lfloor B\rfloor|_S$（命题 5.59）与 $\lfloor B\rfloor$ 的各连通成分既约，存在 $\lfloor B\rfloor$ 的既约成分 $B_1$，使 $B_1\supset D_1\cup D_2$。特别地，在 $D_1\cap D_2$ 的各点处，$X$ 奇异或 $S+B$ 不为 snc。这与 dlt 假定矛盾。
\end{proof}

\begin{corollaryp}{5.62}
设 $X$ 是正规簇，$S$ 是余维 $2$ 且为 Cartier 的既约 Weil 除子，$S'$ 为既约 Weil 除子，$B$ 为有效 $\Q$-除子。若 $(X,S+S'+B)$ 为 dlt，则 $S'|_S$ 的各连通成分既约。
\end{corollaryp}

\begin{proof}
由命题 2.43，可假定 $\lfloor B\rfloor=0$。特别地，$\lfloor S'+B\rfloor$ 既约。因此 $(S,S'|_S+B|_S)$ 为 plt（命题 5.61）。于是，$S'|_S=\lfloor S'|_S\rfloor$ 的各连通成分既约（命题 5.51）。
\end{proof}

% [原文 p.223]
\begin{propositionp}{5.63}
设 $X$ 是正规簇，$S_i$ 是过 $0\in X$ 的 $\Q$-Cartier Weil 除子。若 $\bigl(X,\sum_{i=1}^kS_i\bigr)$ 为 dlt，则 $k\le\dim X$。
\end{propositionp}

\begin{proof}
对 $\dim X$ 归纳证明。$\dim X=1$ 时主张显然。假定 $S_1$ 为 Cartier。设 $T$ 是其既约成分。此时 $T$ 正规（推论 5.52），余维 $2$ 且 Cartier（推论 5.56），$\bigl(T,(S_1-T)|_T+\sum_{i=2}^kS_i|_T\bigr)$ 为 dlt（命题 5.59）。于是由归纳法得 $k\le\dim X$。

其次归约到 $S_1$ 为 Cartier 的情形。设 $m>0$ 是使 $mS_1$ 为 Cartier 的最小正整数。令 $L=\cO_X(S_1)$，$p\colon X'\to X$ 为对应的循环覆盖（定义 2.50）。$p$ 在使 $S_1$ 为 Cartier 的开集合上为 étale。令 $S'_i:=p^{-1}(S_i)$，则 $S'_1$ 为 Cartier。

与命题 5.59 的证明同样选取 $Z\subset X$，令 $Z':=p^{-1}(Z)$。由于 $X\setminus Z$ 非奇异，$(X'\setminus Z')\to(X\setminus Z)$ 为 étale。与命题 5.20 结合，这意味着 $\bigl(X',\sum_{i=1}^kS'_i\bigr)$ 为 dlt。于是如上所证，$k\le\dim X$。
\end{proof}

% [Tutor Note: 命题 5.63 可读作组合上的界：dlt 对在一点处相交的边界成分个数不超过维数。三维时 $k\le3$，这正是 $\Q$-分解 dlt 版 MMP 能逐段推进的理由之一。]
\begin{remark}{5.64}
命题 5.63 是下述结果 [K+92, 18.2] 的特殊情形：若 $(X,\sum b_iS_i)$ 为 lc，则 $\sum b_i\le\dim X$。
\end{remark}

% [原文 p.224]
\section{对偶理论}
\takeaway{本节只处理射影情形，以 CM 层为工具建立所需的对偶理论。核心是 CM 层的 Serre 对偶（定理 5.71）与伴随公式（命题 5.73），二者都无须进入导出范畴，只需层的 $\mathcal{H}om$ 与迹映射。}

本节的目标是解说本书所需的对偶理论结果。Grothendieck 的一般对偶理论（参见 [Har66]）把所有这些结果作为非常特殊的情形包含在内，但 [Har66] 展开的是一般论，从中导出单纯的结果往往很困难。本书所需的结果大多可从 [Har77, III.5-7] 的较为初等的理论导出，但仍需相当多的讨论。

为此，本书决定介绍 CM 层的对偶理论。CM 层的主要优点是，Serre 对偶不必移到导出范畴（derived category）即可起作用。

证明在射影簇的情形非常有效。然而，非射影簇或解析空间的情形须从基础部分起详细讨论。主要问题是，对偶化层的定义（定义 5.66）在非固有簇的情形下没有意义。把命题 5.68 作为一般对偶化层的局部定义采用是可能的。然而，证明这是不依赖紧化的概念并不容易。

因此，只处理射影的情形。一般代数情形让给 [Har66]，复解析情形让给 [BS76]。

为便于引用，只列举应当引用的对偶理论结果。

\par\medskip\noindent\textbf{5.65}（$\PP^n$ 上的 Serre 对偶）\quad 设 $F$ 是（域 $k$ 上的）$\PP^n$ 上的凝聚层。此时，$H^n(\PP^n,F)$ 与 $\operatorname{Hom}_{\PP^n}\bigl(F,\cO_{\PP^n}(K_{\PP^n})\bigr)$ 互为对偶的 $k$-向量空间。若 $F$ 为局部自由，则对每个 $i$，$H^i(\PP^n,F)$ 与 $H^{n-i}\bigl(\PP^n,\cO_{\PP^n}(K_{\PP^n})\otimes F^\vee\bigr)$ 互为对偶的 $k$-向量空间（参见定义 5.2）。这表示为
\[
H^i(\PP^n,F)\overset{\text{对偶}}{\longleftrightarrow}H^{n-i}\bigl(\PP^n,\cO_{\PP^n}(K_{\PP^n})\otimes F^\vee\bigr).
\]
\hfill$\square$

\begin{definition}{5.66}
设 $X$ 是域 $k$ 上的 $n$ 维固有概形。对偶化层（dualizing sheaf）是指凝聚层 $\omega_X$ 与满射 $\operatorname{Trace}_X\colon H^n(X,\omega_X)\to k$ 的组，使得对任意凝聚层 $F$，$\operatorname{Trace}_X$ 给出自然的 $k$-线性同型映射
% [原文 p.225]
\[
\operatorname{Hom}_X(F,\omega_X)\cong\operatorname{Hom}_k\bigl(H^n(X,F),k\bigr)
\]
所定义的对象。$(\omega_X,\operatorname{Trace}_X)$ 若存在，则除自然同型外唯一，这一点容易看出。为简单起见，通常省略 $\operatorname{Trace}_X$。
\end{definition}

% [Tutor Note: 5.65 与定义 5.66 只处理射影情形，因为 $\PP^n$ 上的 Serre 对偶与 $\omega_{\PP^n}$ 已知；一般概形的对偶化层由命题 5.68 经有限态射拉回构造。]
下一个一般结果对于计算对偶化层有用。

\begin{propositionp}{5.67}
设 $f\colon X\to Y$ 是有限态射，$F,G$ 分别是 $X,Y$ 上的凝聚层。令 $f^!G:=\mathcal{H}om_{\cO_Y}(f_*\cO_X,G)$，并在其上自然引入 $\cO_X$-模构造。
\begin{enumerate}[label=(\arabic*),leftmargin=*]
\item 存在自然的 $f_*\cO_X$-线性同型
\[
f_*\mathcal{H}om_{\cO_X}(F,f^!G)=\mathcal{H}om_{\cO_Y}(f_*F,G).
\]
\item 存在自然的 $k$-线性同型
\[
\operatorname{Hom}_X(F,f^!G)=\operatorname{Hom}_Y(f_*F,G).
\]
\end{enumerate}
\end{propositionp}

\begin{proof}
(1) 是局部结果，故可设 $X=\Spec A$，$Y=\Spec B$，$M=\Gamma(X,F)$，$N=\Gamma(Y,G)$。此时，(1) 无非是证明下述等式
\[
\operatorname{Hom}_A\bigl(M,\operatorname{Hom}_B(A,N)\bigr)=\operatorname{Hom}_B(M,N).
\]
其中，左边的元 $\psi$ 与右边的元 $\varphi$ 由 $\psi(m)\colon a\mapsto\varphi(am)$（$a\in A$，$m\in M$）相联系。这给出 $1$ 对 $1$ 对应。

(2) 通过取整体截面由 (1) 得到。
\end{proof}

\begin{propositionp}{5.68}
设 $f\colon X\to Y$ 是纯 $n$ 维固有概形之间的有限态射。若 $\omega_Y$ 存在，则 $\omega_X$ 也存在，且 $\omega_X\cong f^!\omega_Y$ 成立。
\end{propositionp}

\begin{proof}
由命题 5.67(2)，得到自然同型 $\operatorname{Hom}_X(F,f^!\omega_Y)=\operatorname{Hom}_Y(f_*F,\omega_Y)$。由 $Y$ 上的对偶，后者与 $H^n(Y,f_*F)=H^n(X,F)$ 对偶。
\end{proof}

\begin{corollaryp}{5.69}
对域 $k$ 上的任意纯 $n$ 维射影概形 $X$，$\omega_X$ 存在且为 $S_2$。
\end{corollaryp}

\begin{proof}
域 $k$ 上的任意 $n$ 维射影概形有有限态射 $f\colon X\to\PP^n_k$。由 5.65，$\omega_{\PP^n_k}$ 存在。故由命题 5.68，$\omega_X$ 存在。由命题 5.4，$\omega_X$ 为 $S_2$ 当且仅当 $f_*\omega_X$ 为 $S_2$。后者为 $S_2$ 是显然的，因为它是对局部自由层作用 $\mathcal{H}om$ 所得。
\end{proof}

\begin{corollaryp}{5.70}
设 $X$ 是域 $k$ 上的纯 $n$ 维射影概形，$F$ 是 $X$ 上的凝聚层，$\Supp F$ 为纯 $n$ 维。
\begin{enumerate}[label=(\arabic*),leftmargin=*]
\item 若 $F$ 为 CM，则 $\mathcal{H}om_{\cO_X}(F,\omega_X)$ 也为 CM。若 $F$ 为 $S_2$，则逆也成立。
% [原文 p.226]
\item 若 $X$ 为 $S_2$，则 $\cO_X$ 为 CM 当且仅当 $\omega_X$ 为 CM。
\end{enumerate}
\end{corollaryp}

\begin{proof}
设 $f\colon X\to Y\cong\PP^n_k$ 为有限态射。由推论 5.5，$F$ 为 CM 当且仅当 $f_*F$ 局部自由。又，$\mathcal{H}om_{\cO_X}(F,\omega_X)$ 为 CM 当且仅当 $f_*\mathcal{H}om_{\cO_X}(F,\omega_X)=\mathcal{H}om_{\cO_Y}(f_*F,\omega_Y)$ 局部自由。

若 $f_*F$ 局部自由，则 $\mathcal{H}om_{\cO_Y}(f_*F,\omega_Y)$ 局部自由。进一步，若 $f_*F$ 为 $S_2$，则逆也成立（定义 5.2）。此附加条件无非是 $F$ 为 $S_2$（命题 5.4）。因此 (1) 得证。

(2) 是 (1) 的特殊情形。
\end{proof}

\begin{theoremp}[{CM 层的 Serre 对偶（Serre duality）}]{5.71}
设 $X$ 是域 $k$ 上的纯 $n$ 维射影概形，$F$ 是 $X$ 上的 CM 层，$\Supp F$ 为纯 $n$ 维。此时，
\[
H^i(X,F)\overset{\text{对偶}}{\longleftrightarrow}H^{n-i}\bigl(X,\mathcal{H}om_{\cO_X}(F,\omega_X)\bigr).
\]
\end{theoremp}

\begin{proof}
存在有限态射 $f\colon X\to P=\PP^n$，且下式成立。其中注意：应用 Serre 对偶时，$f_*F$ 局部自由（推论 5.5）。
\[
\begin{aligned}
H^i(X,F)&=H^i(P,f_*F) &&(f\ \text{有限})\\
&\overset{\text{对偶}}{\longleftrightarrow}H^{n-i}\bigl(P,\mathcal{H}om_{\cO_P}(f_*F,\omega_P)\bigr)
&&(\PP^n\ \text{上的 Serre 对偶 5.65})\\
&=H^{n-i}\bigl(P,f_*\mathcal{H}om_{\cO_X}(F,\omega_X)\bigr) &&(\text{命题 5.67})\\
&=H^{n-i}\bigl(X,\mathcal{H}om_{\cO_X}(F,\omega_X)\bigr) &&(\text{自明}).
\end{aligned}
\]
\end{proof}

% [Tutor Note: 定理 5.71 的意义在于绕开导出范畴：CM 层的 Serre 对偶只需层 $\mathcal{H}om(F,\omega_X)$，配合有限态射 $f\colon X\to\PP^n$ 与投影公式即可证明。]
最初的归结，是用上同调给出 CM 层的特征刻画。

\begin{corollaryp}{5.72}
设 $X$ 是域 $k$ 上的纯 $n$ 维射影概形，$D$ 为丰富 Cartier 除子。$F$ 是 $X$ 上的凝聚层，$\Supp F$ 为纯 $n$ 维。此时，下述条件等价。
\begin{enumerate}[label=(\arabic*),leftmargin=*]
\item $F$ 为 CM 层。
\item $H^i(X,F(-rD))=0$（$\forall i<n$，$\forall r\gg0$）。
\end{enumerate}
\end{corollaryp}

\begin{proof}
假定 (1)。由定理 5.71，
% [原文 p.227]
\[
H^i(X,F(-rD))\overset{\text{对偶}}{\longleftrightarrow}H^{n-i}\bigl(X,\mathcal{H}om_{\cO_X}(F,\omega_X)(rD)\bigr).
\]
当 $r\gg0$ 且 $i<n$ 时，由 Serre 消没定理，左边 $=0$。于是 (2) 得证。

(2)$\Rightarrow$(1) 对 $n$ 归纳证明。$n=0$ 时无须证明。任取 $x\in X$。由 $H^0(X,F(-rD))=0$，$F$ 不含支集与 $\{x\}$ 一致的部分层。于是，若 $r'\gg0$，则在 $H^0(X,\cO(r'D))$ 的元 $s$ 中，存在在 $x$ 处为 $0$、在 $F$ 的伴随点处不为 $0$ 者。于是 $s\colon F\to F(r'D)$ 为单射。令 $Y:=(s=0)$，得到下述正合列
\[
0\to F\bigl(-(r+r')D\bigr)\to F(-rD)\to F_Y(-rD)\to0.
\]
由 (2) 与正合列，对 $i<n-1$ 与 $r\gg0$ 有 $H^i(Y,F_Y(-rD))=0$。于是由归纳法 $F_Y$ 为 CM。进一步，由命题 5.3(1)，$F$ 在 $x$ 处为 CM。
\end{proof}

\begin{propositionp}[{伴随公式（adjunction formula）}]{5.73}
设 $X$ 是域 $k$ 上的纯 $n$ 维射影 CM 概形，$D\subset X$ 为有效 Cartier 除子，则 $\omega_D\cong\omega_X(D)\otimes\cO_D$ 成立。
\end{propositionp}

\begin{proof}
设 $f\colon X\to P=\PP^n$ 有限，$L\subset P$ 为超平面，$H=f^{-1}(L)$。此处射影空间 $P$、$L$ 的对偶化层已知，故 $\omega_L\cong\omega_P(L)\otimes\cO_L$ 成立。由推论 5.5，$f_*\cO_X$ 局部自由。因此
\[
\begin{aligned}
f_*\omega_H&\cong\mathcal{H}om_{\cO_L}(f_*\cO_H,\omega_L)=\mathcal{H}om_{\cO_P}\bigl(f_*\cO_X,\omega_P(L)\otimes\cO_L\bigr)\\
&=\mathcal{H}om_{\cO_P}(f_*\cO_X,\omega_P)\otimes\cO_P(L)\otimes\cO_L=f_*\bigl(\omega_X(H)\otimes\cO_H\bigr).
\end{aligned}
\]
$D\subset X$ 为任意有效 Cartier 除子时，适当选取有限态射 $f\colon X\to\PP^n$ 与超平面 $L\subset\PP^n$，可使 $H=f^{-1}(L)$ 分解为无公共成分的 $D$ 与 $E$（即 $H=D+E$）。

设 $j\colon D\to H$ 为单射。此时可知 $j_*\omega_D=\mathcal{H}om_{\cO_H}(j_*\cO_D,\omega_H)$ 与 $\omega_H\cong\omega_X(H)\otimes\cO_H$。应证的等式为
\[
\mathcal{H}om_{\cO_X}(\cO_D,\omega_H)\cong\omega_X(D)\otimes\cO_D.
\]
正合列
\[
0\to\cO_X(D)\to\cO_X(H)\to\cO_E(H)\to0
\]
% [原文 p.228]
与 $\omega_X$ 张量积，得到下述正合列
\[
0\to\omega_X(D)\to\omega_X(H)\to\omega_X\otimes\cO_E(H)\to0.
\]
其中，左端的单射性由 $\omega_X$ 为 $S_2$ 得出。进一步，由 $\omega_X$ 为 $S_2$（推论 5.69）可知，$\omega_X\otimes\cO_E(H)$ 作为 $\cO_E$-层为 $S_1$（命题 5.3）。

此处考虑下述正合列
\[
0\to\omega_X(D)/\omega_X\to\omega_X(H)/\omega_X\to\omega_X(H)/\omega_X(D)\to0.
\]
证明了列的正中间项为 $\omega_H$，右边也已算出。改写后得到下述正合列
\[
0\to\omega_X(D)\otimes\cO_D\to\omega_H\to\omega_X\otimes\cO_E(H)\to0.
\]
$\omega_X\otimes\cO_E(H)$ 作为 $\cO_E$-层为 $S_1$，且 $D$ 与 $E$ 无公共成分，故得到下式
\[
\begin{aligned}
j_*\omega_D&=\mathcal{H}om_{\cO_H}(j_*\cO_D,\omega_H)=\mathcal{H}om_{\cO_H}\bigl(j_*\cO_D,\omega_X(D)\otimes\cO_D\bigr)\\
&=\mathcal{H}om_{\cO_D}\bigl(\cO_D,\omega_X(D)\otimes\cO_D\bigr)=\omega_X(D)\otimes\cO_D.
\end{aligned}
\]
\end{proof}

% [Tutor Note: 伴随公式的证明把 $D$ 放进一个超平面截面 $H=D+E$ 中，用 $\omega_H$ 作中转；关键一步是 $\omega_X\otimes\cO_E(H)$ 的 $S_1$ 性，它保证序列可以改写成 $D$ 上的等式。]

\begin{remark}{5.74}
即使 $X$ 不为 CM，由上述证明也可得到单射 $\omega_X(D)\otimes\cO_D\to\omega_D$。然而，这并不限于同型。例如，取 $X$ 为 $\PP^n$ 上的 Abel 簇 $A$ 的锥，顶点为 $0\in X$。$\omega_X$ 是对偶化层。此时 $\omega_X\cong\cO_X$，从而可证 $X$ 在 $0$ 处为 lc。取 $D\subset X$ 为一般超平面截面，则 $D$ 在 $0$ 处不为正则。令 $\bar D$ 为 $D$ 的正规化，则 $\omega_{\bar D}\cong\cO_X(D)|_D\otimes\cO_{\bar D}$ 与 $\omega_X(D)|_D$ 并不同型。$X$ 在 $0$ 处不为 CM，$\omega_X$ 与 [Har66] 的对偶化复形（dualizing complex）不一致。
\end{remark}

% [原文 p.229]
\begin{propositionp}{5.75}
设 $X$ 是域 $k$ 上的 $n$ 维正规射影簇。此时，$\omega_X\cong\cO_X(K_X)$ 成立。
\end{propositionp}

\begin{proof}
此处用 $\omega_{\PP^n}\cong\cO_{\PP^n}(K_{\PP^n})$。把 $X$ 嵌入 $\PP^N$，固定诱导态射 $f\colon X\to P:=\PP^n$ 为有限且函数域扩张可分离的一般射影 $\pi\colon\PP^N\dashrightarrow\PP^n$。此时，表示为 $K_X=f^*K_P+R$。其中 $R$ 的支集在分歧集合上。迹映射（定义 5.6）给出双线性映射 $f_*\cO_X\times f_*\cO_X\to\cO_P$。与 $\cO_P(K_P)$ 取张量积，得到双线性映射 $f_*\cO_X\times f_*\cO_X(f^*K_P)\to\cO_P(K_P)$。它现在
\[
f_*\cO_X\times f_*\bigl(\cO_X(K_X-R)\bigr)\to\cO_P(K_P)
\]
给出。这延拓为
\[
\Phi\colon f_*\cO_X\times f_*\bigl(\cO_X(K_X)\bigr)\to\bigl(\cO_P(K_P)\ \text{的有理截面层}\bigr).
\]
用 $P$ 的局部坐标表示如下。
\[
(g,h\cdot f^*(dx_1\wedge\cdots\wedge dx_n))\mapsto\operatorname{Trace}(gh)dx_1\wedge\cdots\wedge dx_n.
\]
\end{proof}

\par\medskip\noindent\textbf{主张 5.76}\quad $\Phi$ 延拓为（用同样记号表示）下述双线性映射
\[
\Phi\colon f_*\cO_X\times f_*\bigl(\cO_X(K_X)\bigr)\to\cO_P(K_P).
\]
它在使 $f$ 在 $f^{-1}(v)$ 上为平坦的点 $v$ 处非退化。

\begin{proof}
在这样的点处，$f_*\cO_X$ 自由，$f_*(\cO_X(K_X))$ 自反。因此，只需在 $P$ 的余维 $\ge2$ 的部分集合之外证明主张即可。于是，$p\in P$ 是使 $X$ 在 $f^{-1}(p)$ 的所有点处非奇异的点。

此处，分解 $\pi\colon\PP^N\dashrightarrow\PP^n$ 为 $\pi\colon\PP^N\dashrightarrow\PP^{n+1}\dashrightarrow\PP^n$，可使得 $\pi'\colon X\to Y:=\pi'(X)$ 在 $g^{-1}(p)$ 上同型。在 $\PP^{n+1}$ 上取坐标系 $(x_0:\cdots:x_n:y)$，使 $\pi''$ 为到 $x$-坐标的射影，并令 $p=(1:0:\cdots:0)$。以 $(F(x_1,\dots,x_n,y)=0)$ 为 $Y$ 的仿射坐标。$F$ 是关于 $y$ 无重根的多项式，次数 $r=\deg f$。

令
\[
\sigma:=\frac{1}{\partial F/\partial y}f^*(dx_1\wedge\cdots\wedge dx_n),
\]
则 $f_*(\cO_Y(K_Y))$ 的局部自由基由 $y^j\sigma$（$j=0,\dots,r-1$）给出。$1,y,\dots,y^{r-1}$ 给出 $f_*\cO_X$ 的自由基。由引理 5.78，
\[
\operatorname{Trace}(y^j\sigma)=\delta_{j,r-1}\,dx_1\wedge\cdots\wedge dx_n\qquad(j=0,\dots,r-1).
\]
由 $\Phi$ 的定义，$\operatorname{im}\Phi$ 由 $\operatorname{Trace}(y^j\sigma)$（$j=0,\dots,r-1$）生成。因此 $\operatorname{im}\Phi\subset\cO_P(K_P)$。
\[
(y^i,y^j\sigma)=\operatorname{Trace}(y^{i+j}\sigma)=\delta_{i+j,r-1}\,dx_1\wedge\cdots\wedge dx_n\qquad(i+j\le r-1).
\]
因此，内积的行列在副对角线上元素为 $1$，在其上元素为 $0$。
\end{proof}

\begin{proof}[证明（续）]
由于 $f$ 平坦，由主张 5.76 得到
\[
\cO_X(K_X)=\mathcal{H}om_{\cO_P}(f_*\cO_X,\omega_P)=\omega_X
\]
另外，沿 [Har77, III.7.12] 讨论的思路也可给出别的证明。
% [原文 p.230]
由此得证。
\end{proof}

\begin{propositionp}[{相对迹映射（relative trace map）}]{5.77}
设 $X$ 与 $Y$ 是域 $k$ 上的纯 $n$ 维射影概形，$f\colon X\to Y$ 在稠密开集上拟有限。
\begin{enumerate}[label=(\arabic*),leftmargin=*]
\item 存在自然的（不为 $0$ 的）映射 $\operatorname{Trace}_{X/Y}\colon f_*\omega_X\to\omega_Y$。
\item 若 $f$ 双有理且 $X,Y$ 正规，则 $\operatorname{Trace}_{X/Y}$ 在 $f^{-1}$ 为同型的点处同型。
\item 若 $f$ 双有理且 $X,Y$ 非奇异，则 $\operatorname{Trace}_{X/Y}$ 为同型态射。
\end{enumerate}
\end{propositionp}

\begin{proof}
设 $F$ 是 $X$ 上任一凝聚层。此时，对 $j>0$，$\Supp R^jf_*F$ 的维数至多 $n-j-1$。于是，若 $i+j\ge n$ 且 $j>0$，则 $H^i(Y,R^jf_*F)=0$。由 Leray 谱序列，得到下述满射
\[
H^n(Y,f_*F)\to H^n(X,F).
\]
把它用于 $F=\omega_X$，得到下述单射
\[
H^0(X,\cO_X)=H^0\bigl(X,\mathcal{H}om(\omega_X,\omega_X)\bigr)\hookrightarrow H^0\bigl(Y,\mathcal{H}om(f_*\omega_X,\omega_Y)\bigr).
\]
$1\in H^0(X,\cO_X)$ 的像给出映射 $f_*\omega_X\to\omega_Y$。

若 $f$ 双有理，则 $\operatorname{Trace}_{X/Y}$ 为单射。进一步，若 $X$ 与 $Y$ 正规，则 $\operatorname{Trace}_{X/Y}$ 给出不为 $0$ 的映射 $f_*\cO_X(K_X)\to\cO_Y(K_Y)$。$f_*\cO_X(K_X)$ 自然成为 $\cO_Y(f_*(K_X))=\cO_Y(K_Y)$ 的部分层，两者在 $f^{-1}$ 为同型的点处一致。

最后，(3) 是容易的（参见 [Har77, II.8.19]）。
\end{proof}

\begin{lemmap}{5.78}
设 $K$ 为域，$f\in K[x]$ 为无重根的 $r$ 次首一（monic，即最高次系数为 $1$）多项式。设 $V:=K[x]/(f)$ 为 $K$-模。此时，乘以 $\partial f/\partial x$ 的乘法在 $V$ 上可逆，且
\[
\operatorname{Trace}\Bigl(\frac{x^s}{\partial f/\partial x}\Bigr)=\delta_{s,r-1}\qquad(0\le s\le r-1).
\]
\end{lemmap}

\begin{proof}
两个主张都在基扩张下不变。故可假定 $K$ 为代数闭域。表示 $f(x)=\prod_i(x-a_i)$，则 $K[x]/(f)\cong\bigoplus_iK[x]/(x-a_i)$ 成立。此分解的直和成分由 $x$ 与 $\partial f/\partial x$ 的公共特征向量
% [原文 p.231]
生成，故得到下式
\[
\operatorname{Trace}\Bigl(\frac{x^s}{\partial f/\partial x}\Bigr)=\sum_i a_i^s\prod_{j\ne i}\frac{1}{a_i-a_j}.
\]
引理由在 $\Z[A_1,\dots,A_r,\Delta^{-1}]$ 内成立的恒等式导出
\[
g_s(A_1,\dots,A_r):=\sum_iA_i^s\prod_{j\ne i}\frac{1}{A_i-A_j}=\delta_{s,r-1}\qquad(s=0,\dots,r-1).
\]
其中 $\Delta=\prod_{j<i}(A_j-A_i)$。

确认这个恒等式。首先，$g_s$ 关于 $A_i$ 为对称多项式（即对变量置换不变），且 $g_s\cdot\Delta$ 为交代多项式（即两个变量的互换只改变符号），这一点显然。任意交代多项式可被 $\Delta$ 整除。由于 $\deg g_s\cdot\Delta=\deg\Delta-(r-1)+s$，若 $s<r-1$ 则 $g_s=0$，且 $g_{r-1}$ 为常数。为看出 $g_{r-1}=1$，固定 $A_2,\dots,A_r$ 而取极限 $A_1\to\infty$，注意作为取值于 $\Q$ 的函数 $g_{r-1}\to1$。于是得到 $g_{r-1}=1$。
\end{proof}

% [Tutor Note: 引理 5.78 是 Euler 关于部分分式的古典恒等式的特例；恒等式在 $\Z[A_1,\dots,A_r,\Delta^{-1}]$ 内成立，故对任意域 $K$ 都可用。]
\chapterreview{%
本章讨论极小模型纲领运行中必然出现的奇点。第 3、4 章已搭好收缩与翻转的一般框架，然而每一步运算都会产生新的奇异纤维与新的边界，因此必须先精确规定“允许的奇点”指哪一类。本章五节给出这套规定的完整表述。有理奇点刻画结构层与对偶化层的上同调不受奇点消解影响的那些奇点。对数终端给出边界系数为 $1$ 时的局部理论，其主结果是定理 5.22。三维终端奇点与标准奇点借助 Du Val 奇点与椭圆奇点得到离散的分类。伴随反转把边界上的 klt 性或 lc 性反推回原对。对偶理论则以 CM 层为工具，建立 Serre 对偶与伴随公式。

在知识脉络上，本章汇集了三人组的成果。广中的奇点消解与 Grauert–Riemenschneider 的判据给出有理奇点的原型。Elkik、Fujita、Kollár 把有理性推广到对数对的形变。Shokurov 的连通性论证与伴随反转成为归纳的核心工具，[K+92] 第 17、18 章把散落的结果整理成统一形式。与第 2、3 章相比，本章不发明新方法，而是选出恰好够用的类，并证明它们在运算下封闭。

核心结论有五条。有理奇点等价于 $\cO_X\to Rf_*\cO_Y$ 为同型，也等价于对偶化层的高次上同调消没。dlt 是 klt 的对数推广，边界成分正规，且 lc 中心连通（定理 5.48）。三维终端奇点是 cDV 奇点的循环商，指数 $1$ 的分类型式有限，见表 5.1。伴随反转定理 5.50 把 $(S,B|_S)$ 的 klt 性等价为 $(X,S+B)$ 的 plt 性。CM 层的 Serre 对偶（定理 5.71）与伴随公式（命题 5.73）不进入导出范畴，只用一个层的 $\mathcal{H}om$ 就够用。

阅读时须辨析数处。第一，klt 与 dlt 的边界处理不同：dlt 允许边界系数为 $1$，但要求边界正规（命题 5.51）。第二，伴随与伴随反转方向相反，命题 5.46 给出的是伴随方向的不等式，定理 5.50 才是真正的反转，两者不可混用。第三，定义 5.45 的差异系数允许例外除子的中心落在指定闭子集内，故与普通差异系数不同，命题 5.46 中的不等号方向容易记错。第四，定理 5.34 与定理 5.35 的证明集中在编号段 5.36，二者陈述末尾为 $\square$，表示证明另给。第五，表 5.1 的 cA、cAx、cD、cE 是上同调型的记号，不是奇点的最终命名。

本章的叙述有意克制。三维终端奇点的详细分类被作者明确声明本书不需要，只作为背景。逻辑上，本章只用到第 2、3、4 章的工具，但读者若跳过定理 5.48 的连通性证明，就难以理解定理 5.50 与第 6、7 章翻转的归纳结构。此外，本书尾符约定与常见排版相反：证明随陈述之后给出时，陈述末不带符号；证明另给时才印 $\square$。这一约定在本章多次出现，阅读时须留意。
}
